% Ad Atticum 15.3 (ad-atticum-15-03) --- English translation
% Translated by Claude (Anthropic) from the Latin (Perseus).
% Style guide: STYLE.md.

\ciceroLetterOpener{CICERO ATTICO salutem}

\ciceroSection{1} On the eleventh before the Kalends I received in the Arpinate two letters from you, in which you replied to two of mine. One was dated the fifteenth, the other the twelfth before the Kalends. To the earlier, then, first. You will hurry to the Tusculan villa, as you write; where I had thought I should arrive on the sixth before the Kalends. As to your saying that one must obey the victors --- not I, at any rate, who have many things weightier. For those proceedings you call to mind, conducted in the temple of Apollo in the consulship of Lentulus and Marcellus, the case is not the same nor is the season comparable, especially when you write that Marcellus and others are taking themselves off. We must, then, scent it out face to face and settle whether we can be safely at Rome. The settlers in the new colony certainly do move me; for we are turning amid great straits. But those things are small; indeed we hold even greater in contempt. I have learned the will of Calva, a base and sordid man. The will-tablet of Demonicus, which you care for, I am grateful for. About \dots\ I have already long since written to Dolabella in the fullest detail, provided the letter has been delivered. On his account I both wish and owe it.

\ciceroSection{2} I come to the nearer letter. About Alexio I have learned what I wanted. Hirtius is on your side. As for Antony, since he is what he is, I want it the worse for him. About young Quintus, as you write, \dots; about the father we shall deal face to face. Brutus I desire to help with every means I can. As for his little speech, I see that you feel about it as I do. But I scarcely understand what you want me to write --- as though it were a speech delivered by Brutus, when he himself has published it. How, after all, does that fit together? Or in the way I should write of a tyrant most justly slain? Many things will be said, many written by me, but in another manner and at another time. About Caesar's chair --- well done by the tribunes; splendid too about the fourteen rows! That Brutus was with me, I am glad --- provided he was there willingly and long enough.
